\section{Groups with loops and their rigidity}
\label{sec:part-a}
We use $[x,y]=xyx^{-1}y^{-1}$. A peripheral subgroup is indexed even when two indices give the same subgroup. The notation $\Out_{\mathrm{lac}}$ means preservation of the indexed peripheral conjugacy classes, with their permitted permutations; in hyperbolic type these permutations are determined by the automorphism. We specify the additional data needed in the exceptional types below. All profinite homomorphisms and automorphisms in this section are continuous.

\begin{claim}{A.1}{Proved}{The presentation, ranks, and exceptional types}
Let $X$ be a connected closed oriented topological surface of genus $g$, let $S$ consist of $\nu$ distinct points, and put $U=X\setminus S$. A discrete group with loops of type $(g,\nu)$ consists of a group $\pi$, a finite index set $I$ of cardinality $\nu$, a free rank-one integral orientation module $T$, and exterior homomorphisms $\kappa_i:T\to\pi$. Here an exterior homomorphism is a conjugacy class, under conjugation in the target, of actual homomorphisms. After choosing a generator of $T$, representatives $c_i=\kappa_i(1)$, an ordering of $I$, and the handle generators, these data admit the presentation
\[
 \pi_{g,\nu}=\left\langle a_1,b_1,\ldots,a_g,b_g,c_1,\ldots,c_\nu\ \middle|\
 \prod_{j=1}^{g}[a_j,b_j]c_1\cdots c_\nu=1\right\rangle,
 \qquad L_i=\langle c_i\rangle.
\]
The profinite version uses a free rank-one $\widehat\Z$-module $T$, continuous exterior homomorphisms, and the corresponding presentation in profinite groups, namely the quotient of the free profinite group by the closed normal subgroup generated by this relator. When $I$ is empty one additionally identifies $T$ with the dual of $H^2(\pi;\Z)$, or its continuous profinite counterpart. This is the covariant orientation convention of scan 116. More precisely, an earlier cohomological marking $\kappa:H^2(\pi;\Z)\xrightarrow{\sim}T_{\mathrm{coh}}$ on scan 91 gives, by dualization, $\kappa^\vee:T_{\mathrm{coh}}^\vee\xrightarrow{\sim}H_2(\pi;\Z)$; the covariant module used here is $T=T_{\mathrm{coh}}^\vee$. The same typed dual map specifies the profinite convention. For an isomorphism of groups with loops the data are an isomorphism $f:\pi\to\pi_1$, a bijection $s:I\to I_1$, and an isomorphism $t:T\to T_1$ such that $f\kappa_i$ is conjugate to $\kappa_{s(i)}t$ for every $i$, with the corresponding dual-cohomology compatibility when $I$ is empty. An exterior group with loops regards two such $f$ as equivalent when they differ by target inner conjugation, with the same $s,t$. Thus its isomorphism sets are explicitly the quotient of the preceding isomorphism sets by $\Inn(\pi_1)$. In all nontrivial admitted types $L_i=\kappa_i(T)$ recovers its source by A.2; the one-puncture sphere requires the original source map.
For $\nu>0$ the group is free of rank $2g+\nu-1$, and its abelianization has that rank. For $\nu=0$ the abelianization has rank $2g$. The abelian types are exactly
\[
 (0,0),\quad(0,1),\quad(0,2),\quad(1,0),
\]
with groups $1,1,\Z,\Z^2$, respectively. The same list holds for their profinite completions. The peripheral maps from the source orientation group $\Z$ to $\pi$ are injective except at $(0,1)$, among the types admitted in scans 91--92. The closed sphere is explicitly excluded there.
\end{claim}
\begin{proof}
Eliminate $c_\nu$ when $\nu>0$; the remaining generators have no relation. Abelianizing the displayed presentation leaves only $c_1+\cdots+c_\nu=0$, giving both ranks. A free group of rank at least two is nonabelian, as witnessed by its quotient $S_3$, generated by $(12)$ and $(123)$. A closed genus-$g$ group for $g\ge2$ maps onto the free group on $a_1,\ldots,a_g$ by killing all the $b_j$. The same finite quotient witnesses noncommutativity after completion. For $\nu\ge2$, each peripheral generator belongs to a free basis obtained by eliminating a different peripheral generator. For $\nu=1$, $g\ge1$, the nonempty cyclically reduced word $\prod[a_j,b_j]$ is nontrivial and has infinite order. Its image has infinite order in every pro-$p$ completion: map the first pair into the upper unitriangular $3\times3$ matrices over $\Z_p$, and all other pairs to $1$; the commutator is $1+E_{13}$. This also proves that its procyclic closure in the full completion is $\widehat\Z$, since every prime component injects. At $(0,1)$ the peripheral word is $1$, so the source map $\Z\to1$ is not injective. These calculations prove the correction $\nu\le2$ to the bracket $\nu\le3$ on scan 116. This is mathematical evidence of an error, not a changed reading of the visible numeral.
\end{proof}
\noindent\emph{Manuscript:} \S12, folios 44--48, scans 91--100; \S13 bis, folios 56--57, scans 115--118.

\begin{claim}{A.2}{Proved}{Canonical comparison of the orientation modules}
Assume $\nu>0$ and $\pi\ne1$. There is a canonical exact sequence
\[
0\longrightarrow T_\pi\longrightarrow\bigoplus_{i\in S}L_i
 \longrightarrow\pi^{\mathrm{ab}}\longrightarrow H_1(X;\Z)\longrightarrow0,
\qquad
T_\pi=\ker\left(\bigoplus_iL_i\longrightarrow\pi^{\mathrm{ab}}\right).
\]
Each projection $T_\pi\to L_i$ is an isomorphism of rank-one integral modules. If $M\subset U$ is obtained by removing small open discs at the punctures, the boundary morphism and evaluation give canonical isomorphisms
\[
 H_2(M,\partial M;\Z)\xrightarrow{\sim}T_\pi,
 \qquad
 T_\pi\xrightarrow{\sim}\Hom\bigl(H^2_c(U;\Z),\Z\bigr).
\]
For $\nu=0$, $g\ge1$, the corresponding formula is
$T_\pi=\Hom(H^2(\pi;\Z),\Z)=H_2(X;\Z)$.
Filling the discs identifies both orientation lines with $H_2(X;\Z)$, naturally for homeomorphisms of the marked surface. All these comparisons persist with $\widehat\Z$ and continuous relative group cohomology. At $(0,1)$ the kernel formed from the image subgroup $L_1=1$ is zero, although $H^2_c(\R^2;\Z)=\Z$; the source peripheral morphism is therefore essential. At $(0,0)$ group cohomology $H^2(1;\Z)=0$ cannot supply the surface orientation line.
\end{claim}
\begin{proof}
The relevant portion of the integral homology sequence of the compact pair is
\[
0\to H_2(M,\partial M)\xrightarrow{\partial}
H_1(\partial M)\to H_1(M)\to H_1(M,\partial M).
\]
The initial zero follows from $H_2(M)=0$. A relative fundamental class has boundary the sum of the consistently oriented boundary circles. Identifying each boundary circle with its peripheral subgroup, the boundary morphism identifies $H_2(M,\partial M)$ with the stated kernel and has degree $1$ on each factor. The cokernel is $H_1(X)$ by attaching the discs and using the homology sequence. Universal coefficients for the pair, whose homology is free, identify $H^2(M,\partial M)$ with the dual of $H_2(M,\partial M)$; excision identifies the former with $H^2_c(U)$. Let $D\subset X$ be the union of the filled discs. The canonical isomorphisms $H_2(X)\xrightarrow{\sim}H_2(X,D)\xrightarrow{\sim}H_2(M,\partial M)$ are, respectively, the pair map and excision. Their inverse is the filling comparison, sending a relative fundamental class to the closed fundamental class and proving compatibility without a chosen sign. Conjugating a peripheral subgroup changes its identification by an inner normalizer action, which is trivial by A.3; hence the sequence is independent of paths to the boundary. For the annulus the same conclusion follows directly because its group is abelian.

For clarity, the profinite compact-support notation is the relative continuous cochain complex
\[
 C^*_{\mathrm{rel}}(\widehat\pi;(\widehat L_i);\Z/n)
 =\operatorname{Cone}\left(C^*_{\mathrm{cont}}(\widehat\pi;\Z/n)
 \longrightarrow\bigoplus_iC^*_{\mathrm{cont}}(\widehat L_i;\Z/n)\right)[-1].
\]
For $\nu>0$, $\widehat\pi$ has cohomological dimension one and each $\widehat L_i\simeq\widehat\Z$; thus its degree-two group is the cokernel of restriction on $H^1$. In the displayed presentation the restriction matrix is the transpose of the integral peripheral inclusion matrix. Its cokernel is canonically $\Hom(T_\pi,\Z/n)$. Taking the inverse limit over $n$ gives the claimed duality. For a closed surface of genus at least two, the continuous duality pairing identifies $H^2(\widehat\pi;\widehat\Z)$ with the completed surface orientation cohomology; this is the surface Poincare duality statement in \cite[Proposition 1.3]{A_Mochizuki2007}. At genus one, for each prime $p$ the Koszul resolution of the trivial module over $\Z_p[[\Z_p^2]]$ has differentials given by the two generators minus $1$. Applying trivial coefficients $\Z/p^r$ makes both differentials zero, so the cup product identifies degree two with the exterior square of degree one, exactly as for the discrete torus. The other prime factors have zero positive-degree cohomology with these coefficients. The Chinese remainder theorem gives the comparison for every finite coefficient $\Z/n$. The transition maps on degree one are surjective, so their first derived inverse limit vanishes; consequently passing to inverse limits introduces no extra degree-two term. The canonical map from the discrete presentation to its completion induces the identity on the surface relation and on every peripheral source; consequently the finite-coefficient comparisons commute with filling. The two trivial-group exceptions follow by the explicit cohomology calculations in the statement.
\end{proof}
\noindent\emph{Manuscript:} \S13 bis, folios 56--57, scans 116--118; \S14, folio 64 verso, scan 130. We read $H^2_!$ on scan 130 as the relative complex just specified, rather than as ordinary cohomology of the underlying abstract profinite group.

\begin{claim}{A.3}{Proved}{Discrete peripheral rigidity in every hyperbolic type}
For $2g+\nu>2$, the group $\pi=\pi_{g,\nu}$ has trivial centre and $C_\pi(H)=1$ for every finite-index subgroup $H$. For every peripheral subgroup and every integer $n\ge1$,
\[
 N_\pi(L_i)=C_\pi(L_i)=C_\pi(L_i^n)=L_i.
\]
Distinct peripheral indices give distinct, noncommensurable conjugacy classes.
\end{claim}
\begin{proof}
For $\nu>0$, the free action of a free group on its Cayley tree shows that a nonidentity element has an axis and that its centralizer is the infinite cyclic group of translations generated by its primitive root. An element normalizing this cyclic group preserves the axis. A reversal of the axis fixes a vertex or an edge midpoint, which is incompatible with the free, inversion-free action; the normalizer is the same cyclic group. For $\nu\ge2$ each peripheral element is primitive in a free basis. For $\nu=1$, the cyclic word $\prod[a_j,b_j]$ contains the positive letter $a_1$ exactly once. A proper cyclic power has the number of occurrences of every signed letter divisible by its exponent, so this word is not a proper power. This proves the peripheral formulas including the marginal case of scan 117.

For the closed case, the hyperbolic polygon realization of \cite[Theorem 1.2]{A_FarbMargalit2012} embeds $\pi_g$ as a torsion-free cocompact group of orientation-preserving isometries of the hyperbolic plane. Every nonidentity element is hyperbolic. Its centralizer preserves its axis and acts discretely by translations, hence is cyclic. An axis reversal would be an order-two elliptic element and is absent. A finite-index subgroup is a surface group of genus $1+[\pi:H](g-1)\ge2$. In the punctured case the covering has free rank $1+[\pi:H](2g+\nu-2)\ge2$. Thus a finite-index subgroup cannot lie in a cyclic centralizer. This proves $C_\pi(H)=1$, including $H=\pi$.

For distinct punctures with $\nu\ge3$, their primitive abelianization vectors span distinct lines; for $\nu=2$, $g\ge1$, choose the free basis containing $c_1$. The cyclically reduced lengths of $c_1$ and $c_2=(\prod[a_j,b_j]c_1)^{-1}$ are $1$ and $4g+1$, and their abelianizations force any commensurability to match equal absolute exponents. They cannot be conjugate. This proves the final assertion; at $\nu=1$ it is vacuous.
\end{proof}
\noindent\emph{Manuscript:} \S12, folios 47 verso--48, scans 98--99; \S13 bis, folio 57, scans 117--118.

\begin{claim}{A.4}{Proved}{Exact exceptional-case table}
The discrete and profinite exceptional groups have the following properties; in the profinite column replace $\Z$ by $\widehat\Z$.
\[
\begin{array}{c|c|c|c|c}
(g,\nu)&\pi& Z(\pi)&C_\pi(H)\ (H\text{ finite index/open})&\text{peripherals}\\\hline
(0,0)&1&1&1&\text{none}\\
(0,1)&1&1&1&L_1=1\\
(0,2)&\Z&\Z&\Z&L_1=L_2=\pi\\
(1,0)&\Z^2&\Z^2&\Z^2&\text{none}
\end{array}
\]
Consequently the peripheral normalizer and power-centralizer equalities hold in the plane and annulus, while distinctness of the two peripheral classes fails on the annulus. The nonzero vector $1$ on the annulus and $(1,0)$ on the torus witness failure of both centre and finite-index-centralizer assertions. The plane's failure is the noninjective source map, not a nontrivial centre. For the standard base-element formula of scan 501, the annulus satisfies the equality since its basis element generates its whole group. On the torus $C(\langle(1,0)\rangle^n)=\Z^2$ for every $n\ge1$, and $(0,1)$ witnesses its failure to equal $\langle(1,0)\rangle$. The same vector works after completion. The sphere and plane have no nontrivial base elements to test.
\end{claim}
\begin{proof}
A.1 identifies the groups and peripheral images. Every subgroup of an abelian group is centralized by the entire group. The assertions now follow by direct substitution. For the profinite version the listed elements remain nonzero, as their reductions modulo $2$ show.
\end{proof}
\noindent\emph{Manuscript:} \S20, folios 111--115 verso, scans 222--230; omissions on scans 98 and 116.

\begin{claim}{A.5}{Proved}{Peripheral decomposition groups and transporters}
Let $1\to\pi\to E\xrightarrow{p}\Gamma\to1$ arise from an action on a hyperbolic marked surface. Choose one peripheral subgroup $L_s$ and the corresponding peripheral lift group $E_s\subset E$; thus $E_s\cap\pi=L_s$ and $p(E_s)=\Gamma_s$, the stabilizer of that puncture. Define
\[
\operatorname{Trans}_E(L_s^d,L_s)
 =\{e\in E:eL_s^de^{-1}\subseteq L_s\}\qquad(d\ge1).
\]
Then $N_E(L_s)=\operatorname{Trans}_E(L_s^d,L_s)=E_s$. The same assertions hold for the corresponding profinite extension and closed peripheral groups. On the annulus these assertions fail for $z\mapsto1/z$: with $E=\Z\rtimes C_2$ and the involution acting by inversion, $E_s=\Z$ but $N_E(L_s)=E$ and every transporter is $E$. More generally for every annulus action both normalizer and transporter equal $E$, while $E_s$ is the inverse image of the end stabilizer. Thus their equality with $E_s$ holds exactly when the action fixes both ends. On the one-puncture sphere the peripheral mixed group maps onto $E=\Gamma$ with kernel $\Z$; using its image, both equalities hold, whereas identifying that mixed group itself with a subgroup is invalid. The closed sphere and torus have no peripheral statement.
\end{claim}
\begin{proof}
A normalizer element maps to $\Gamma_s$, since the peripheral conjugacy classes are distinct. Choose $e_s\in E_s$ with the same image. Then $e_s^{-1}e$ belongs to $N_\pi(L_s)=L_s$, proving the normalizer equality. For the transporter assertion, $eL_s^de^{-1}\subset L_s$ implies $L_s^d\subset e^{-1}L_se$. The latter is abelian, so it lies in $C_\pi(L_s^d)=L_s$, by A.3 or A.7. In the discrete case a peripheral generator cannot be a proper power, so the inclusion of peripheral groups $e^{-1}L_se\subset L_s$ is an equality. In the profinite case use compactness of $E$: a closed subgroup $L$ of a compact group cannot be properly contained in an inner conjugate. Indeed $aLa^{-1}\subset L$ gives $a^kLa^{-k}\subset L$ for all $k\ge1$, and a net of positive powers of $a$ converges to $a^{-1}$ inside its procyclic closure. Passing to the limit gives the reverse conjugation inclusion and hence equality. Thus the transporter normalizes $L_s$ in both cases. This argument only uses the exact power-centralizer statement and avoids replacing the one-sided edge statement in \cite[Proposition 1.2(i)]{A_Mochizuki2007} by a stronger unstated assertion about chosen subgroup representatives. The peripheral lift groups for the action on a punctured disc map onto the puncture stabilizer by construction, so the preceding choice of $e_s$ is legitimate. For the annulus the involution swaps the two punctures, whose inertia subgroups both equal $\pi$. Its lift is the explicit element in the normalizer outside $E_s$. Replacing $\Z$ with $\widehat\Z$ gives the completed counterexample.
\end{proof}
\noindent\emph{Manuscript:} \S9, folio 33, scan 69; \S12, folio 48, scan 99. In the plane, the map from the peripheral mixed group to $E$ has a kernel $\Z$; writing that mixed group as a subgroup would already discard essential data.

\begin{claim}{A.6}{Proved}{Centre-free extension reconstruction and its abelian correction}
For a centre-free discrete group $\pi$ and a homomorphism $\rho:\Gamma\to\Out(\pi)$, the extension inducing $\rho$ is, up to a unique isomorphism inducing the identity on its prescribed kernel and quotient,
\[
 E_\rho=\{(a,\gamma)\in\Aut(\pi)\times\Gamma:[a]=\rho(\gamma)\}.
\]
The same formula holds for a topologically finitely generated centre-free profinite $\pi$ and continuous $\rho$, with the profinite topology. For $A=\Z$ or $\Z^2$ and a prescribed action $\rho:\Gamma\to\Aut(A)$, extension classes with their kernel and quotient identified are instead classified by $H^2(\Gamma;A_\rho)$. Their automorphisms inducing the identity on both ends form $Z^1(\Gamma;A_\rho)$. For finite $\Gamma$ this statement holds with $A=\widehat\Z$ or $\widehat\Z^2$ and continuous cochains. Thus both the annulus and the torus require the abelian correction.
\end{claim}
\begin{proof}
An extension $E$ maps to $E_\rho$ by $e\mapsto(\operatorname{conj}_e,p(e))$. Its kernel is $Z(\pi)=1$. Given $(a,\gamma)$, choose $e_0$ above $\gamma$. The automorphism $a\operatorname{conj}_{e_0}^{-1}$ is inner, implemented by a unique element of $\pi$, so the map is onto. An isomorphism inducing the identity on both ends fixes both components of this formula and is unique. In the profinite case this continuous bijection between compact Hausdorff groups is an isomorphism.

For abelian $A$, choose a set-theoretic section $t$ of the quotient with $t(1)=1$ and write $t(g)t(h)=c(g,h)t(gh)$. Associativity is exactly
$c(g,h)+c(gh,k)=\rho(g)c(h,k)+c(g,hk)$. Changing $t$ by a function $b:\Gamma\to A$ adds its coboundary. Conversely this identity makes $(a,g)(b,h)=(a+\rho(g)b+c(g,h),gh)$ associative. An automorphism fixing both ends has the form $(a,g)\mapsto(a+b(g),g)$, where the homomorphism equation says $b$ is a $1$-cocycle. For finite $\Gamma$ all choices of section are continuous. For example $\Gamma=C_2$, $A=\Z$ with trivial action has the two extension classes $\Z\times C_2$ and $\Z$, with kernel $\Z\xrightarrow{\times2}\Z$ in the latter. They induce the same outer action and are not isomorphic as extensions; reduction modulo $2$ detects the different value of $t(\sigma)^2$.
\end{proof}
\noindent\emph{Manuscript:} \S12, folio 47 verso, scan 98; \S17, folios 77--80, scans 155--160.

\begin{claim}{A.7}{Known}{The full profinite rigidity statements, with reductions}
For the full profinite completion $\Pi=\widehat\pi_{g,\nu}$ of a hyperbolic type, $\Pi$ is torsion-free, $C_\Pi(H)=1$ for every open subgroup $H$, and
\[
N_\Pi(\widehat L_i)=C_\Pi(\widehat L_i^n)=\widehat L_i
\quad(n\in\Z_{>0}).
\]
Distinct peripheral classes are noncommensurable. The same conclusions hold for the pro-$p$ completions, for every prime $p$, including $2$ and $3$. The four exceptional types have exactly the completed table of A.4, and A.5 supplies the extension and transporter assertions.
\end{claim}
\begin{proof}
Take the full formation of finite groups, respectively the formation of finite $p$-groups, in \cite[Remark 1.2.2 and Proposition 1.4]{A_MochizukiTamagawa2008}. A hyperbolic complex curve with the prescribed topological presentation identifies its geometric group with $\Pi$, respectively the pro-$p$ completion. Proposition 1.4 asserts precisely the vanishing of centralizers of open subgroups, with no properness assumption, so taking $H=\Pi$ gives centre-freeness. Remark 1.2.2 gives finite $\ell$-cohomological dimension for every relevant prime, hence excludes an element of order $\ell$: such an element would give a closed $C_\ell$ subgroup, whose $\ell$-cohomological dimension is infinite. Every nontrivial finite-order element has a prime-order power; this proves torsion-freeness.

For peripherals, \cite[Lemma 1.3.7]{A_Mochizuki2004} gives $\operatorname{Comm}_\Pi(\widehat L_i)=\widehat L_i$. A normalizer is contained in that commensurator. An element centralizing $\widehat L_i^n$ conjugates the finite-index subgroup $n\widehat L_i$ to itself, so it too belongs to the commensurator. Conversely the procyclic group $\widehat L_i$ centralizes every power. This proves both equalities. In the semi-graph of the smooth pointed curve, peripherals are the cusp edge groups. Proposition 1.2(i) of \cite{A_Mochizuki2007} says that an intersection open in one edge group identifies the underlying edge; hence distinct punctures cannot have commensurable inertia groups. This gives the needed distinction of peripheral indices; A.5 proves its transporter assertion separately using power-centralizers and compactness. No claim about centralizers of arbitrary non-peripheral elements is needed here.

For $\nu\ge2$ there is also a direct source for the power-centralizer assertion: each peripheral generator can be made a free basis element, and \cite[Lemma (2.1.2)]{A_Nakamura1994} gives the centralizer of every nonidentity element of its procyclic closure. For $\nu=1$ the commensurator argument covers precisely the case that this basis argument does not cover. For every $\nu\ge1$, each handle generator $a_j,b_j$ belongs to the free basis obtained by eliminating $c_\nu$. Applying that same lemma proves equation (4) of scan 501 for all these base elements, including the once-punctured case. The basis condition is material if the phrase ``systeme de generateurs on scan 500 is read literally without a cardinality restriction: in a free profinite group on $a,b$, the redundant generating list $a,b,a^2$ has $C(a^2)=\overline{\langle a\rangle}\supsetneq\overline{\langle a^2\rangle}$. The last index is $2$, detected in abelianization modulo $2$. For the pro-$p$ version use $a,b,a^p$ and index $p$. The stated source theorem applies to a member of a free basis; it does not supply the assertion for every redundant generating list.
\end{proof}
\noindent\emph{Manuscript:} \S13 bis, folio 57 verso, scan 118; \S35, folios 253--254, scans 499--501.

\begin{claim}{A.8}{Proved}{Lemme 5 from open-subgroup centralizers alone}
Let $P$ be a discrete or profinite group for which the centralizer of every finite-index, respectively open, subgroup is trivial. Let $H\subset P$ be such a subgroup and let $u\in\Aut(P)$ preserve $H$. An automorphism restricting to the identity on $H$ is the identity on $P$. An automorphism whose restriction equals $\operatorname{conj}_h$ for some $h\in H$ equals $\operatorname{conj}_h$ on $P$.
\end{claim}
\begin{proof}
Take the finite-index normal core $N=\bigcap_{x\in P}xHx^{-1}\subset H$. If $u|_H=\id$, then for $x\in P$ and $n\in N$,
\[
 u(x)nu(x)^{-1}=u(xnx^{-1})=xnx^{-1}.
\]
Thus $x^{-1}u(x)\in C_P(N)=1$ for every $x$. For the second assertion apply the first to $\operatorname{conj}_{h^{-1}}u$. In the profinite case the core is an intersection of finitely many open subgroups and is open. This proof uses neither peripheral generators nor their power-centralizers.
\end{proof}
\noindent\emph{Manuscript:} \S35, folio 252 verso, scan 498, Lemme 5; folios 253--254, scans 500--501, proposed proof.

\begin{claim}{A.9}{Proved}{Faithfulness ascends along a finite etale cover}
Let $k$ be a characteristic-zero field, $U/k$ a hyperbolic geometrically connected smooth curve, and $V\to U$ a finite etale geometrically connected covering defined over $k$. Put $\Gamma=\operatorname{Gal}(\overline k/k)$ and choose compatible geometric base points. Faithfulness of $\Gamma\to\Out(\pi_1(U_{\overline k}))$ implies faithfulness of $\Gamma\to\Out(\pi_1(V_{\overline k}))$.
\end{claim}
\begin{proof}
The covering gives an open embedding $\Pi_V\hookrightarrow\Pi_U$ of arithmetic fundamental groups, compatible with their surjections onto the same $\Gamma$. Their geometric kernels are an open subgroup $H\subset P$. Let $\gamma\in\Gamma$ act trivially in $\Out(H)$ and choose a lift $e\in\Pi_V$. Conjugation by $e$ preserves $H$ and restricts there to conjugation by some $h\in H$. By A.7 and A.8 its action on $P$ is also conjugation by $h$. Thus $\gamma$ acts trivially in $\Out(P)$ and equals $1$. A change of compatible base points conjugates these homomorphisms and does not change their kernels. The proof does not use the structure of finite subgroups of $\Gamma$.
\end{proof}
\noindent\emph{Manuscript:} \S35, folio 253, scans 499--500, Corollaire 2.

\begin{claim}{A.10}{Known}{Closed-surface base elements: two separate source reductions}
Let $\Pi=\widehat\pi_g$ with $g\ge2$, let $x$ be any of the standard surface generators $a_i,b_i$, and let $\Lambda=\overline{\langle x\rangle}$. Then $\Lambda\simeq\widehat\Z$ and, for every positive integer $n$,
\[
 C_\Pi(\Lambda^n)=N_\Pi(\Lambda^n)=\Lambda.
\]
\end{claim}
\begin{proof}
The standard generator maps to a primitive vector of $\Pi^{\mathrm{ab}}=\widehat\Z^{2g}$; hence its procyclic closure injects there and is $\widehat\Z$. For the first source reduction, $x$ is represented by a nonseparating simple closed curve, and is therefore algebraically simple in the terminology of \cite[Definition 1.3]{A_BoggiZalesskii2019}. Proposition 3.6 of that paper gives $N_\Pi(\Lambda^n)=\Lambda$. Centralizers lie in normalizers and contain $\Lambda$, proving the centralizer equality.

For the second reduction, pinch that curve on the genus-$g$ surface. The resulting stable one-node curve has a normalization of genus $g-1$ with two marked inverse branches; at $g=2$ this is a twice-marked genus-one component, still hyperbolic. Its associated semi-graph of anabelioids has one edge group whose image in the smoothing's profinite group is exactly $\Lambda$, up to the chosen base-path conjugation. Proposition 1.2(ii) of \cite{A_Mochizuki2007} makes this edge group commensurably terminal. An element normalizing $\Lambda^n$ commensurates $\Lambda$, proving the same normalizer equality. These are independent applications of the two stated source results, and neither replaces a closed surface group by a free group.
\end{proof}
\noindent\emph{Manuscript:} \S35, folio 254, scan 501, equation (4).

\begin{claim}{A.11}{Proved}{What the abstract profinite presentation specifies}
The full profinite presentation required by scan 318, equation (4), already determines a profinite group with loops isomorphic to $\widehat\pi_{g,\nu}$ with its standard loops. Hence A.7 and A.10 apply to every abstract object satisfying that presentation axiom. Merely requiring generators to satisfy the displayed equation would not suffice: $P=\widehat\Z^2$, $c_1=(1,0)$, $c_2=(0,1)$, $c_3=(-1,-1)$ satisfies the $(0,3)$ relation, has three injective procyclic peripheral subgroups and the rank-one orientation kernel, but has centre $P$.
\end{claim}
\begin{proof}
The universal property gives a continuous homomorphism from the quotient of the free profinite group by the closed normal closure of the relator to $P$. The assertion that this is a defining relation makes this homomorphism an isomorphism; its inverse sends the specified base to the standard base. This gives exactly the comparison map in scan 319, equation (5). For the weaker condition the explicit abelian example has $c_1+c_2+c_3=0$ and generates $P$, but has the additional relation $[c_1,c_2]=1$. Its peripheral inclusion matrix has kernel generated by $(1,1,1)$, and its centre contains $(1,0)\ne0$. The finite reductions and the rank calculation are checked by the accompanying code.
\end{proof}
\noindent\emph{Manuscript:} \S26, folio 161, scans 318--319, equations (4)--(5). The words ``relation generatrice'' are indispensable; weakening them changes the objects by the explicit quotient above.

\begin{claim}{A.12}{Proved}{The discrete filling isomorphism and all compatibility maps}
Let $U=X\setminus S$ be hyperbolic, choose $b\in U$, and put $U'=U\setminus\{b\}$, $\pi'=\pi_1(U')$, and $\pi=\pi'/\langle\!\langle L_b\rangle\!\rangle$. Choose the peripheral representative $L_b\subset\pi'$. The normalization of representatives defines an isomorphism of discrete groups
\[
 d_b:\Out_{\mathrm{lac}}(\pi';b)\xrightarrow{\sim}\Aut_{\mathrm{lac}}(\pi).
\]
On both sides old punctures may be permuted and orientation may be reversed. It identifies the pure orientation-preserving subgroups, whose common index is $2\nu!$, where $\nu=|S|$. It commutes with forgetting the distinguished puncture, with the maps to $\Out_{\mathrm{lac}}(\pi)$, with orientation characters, with the permutation maps to $\mathfrak S_\nu$, and with point-pushing and inner conjugation on their kernels.
\end{claim}
\begin{proof}
An outer class fixing $[L_b]$ has a representative $\alpha$ preserving $L_b$ itself: conjugate any representative by an element taking its peripheral image back to $L_b$. Two normalized representatives differ by conjugation by an element of $N_{\pi'}(L_b)=L_b$, by A.3. Such conjugation becomes the identity after killing $L_b$. Thus $\alpha$ induces a well-defined actual automorphism $d_b([\alpha])$ of the specified quotient $\pi$. Products of normalized representatives are normalized, so $d_b$ is a homomorphism. Its definition also proves compatibility with the permutation of the remaining loops and with the orientation module, using A.2's filling comparison.

We give the group-theoretic identification of this homomorphism with the topological one. The Dehn--Nielsen--Baer theorem identifies the extended mapping class group of a hyperbolic punctured surface with the outer automorphisms preserving its peripheral conjugacy classes, and in the closed case with $\Out(\pi)$; the precise versions are \cite[Theorems 8.1 and 8.8]{A_FarbMargalit2012}. The evaluation map of the homeomorphism group of $(X,S)$ at $b$ is a locally trivial fibration over $U$: over a sufficiently small disc, move $b$ to a nearby point by a homeomorphism supported in that disc, which constructs a local section. Its connecting map pushes $b$ round a loop. The identity component of this homeomorphism group is contractible in hyperbolic type, by the surface homeomorphism-group theorem in \cite[Lemma 2.1 and Theorem 1.1]{A_Yagasaki2000}; hence the homotopy sequence gives
\[
1\to\pi\xrightarrow{\mathrm{push}}\Out_{\mathrm{lac}}(\pi';b)
\xrightarrow{\mathrm{fill}}\Out_{\mathrm{lac}}(\pi)\to1.
\]
One may alternatively use the extended Birman exact sequence, \cite[Theorem 4.6 and the immediately following extended-group remark]{A_FarbMargalit2012}. Surjectivity follows directly by moving the image of $b$ back to $b$ along an arc avoiding $S$.

Choose the convention that the point-push associated to a based loop $\gamma$ moves the base-point track in the direction for which the induced based map is $\eta\mapsto\gamma\eta\gamma^{-1}$. This can also be checked on the boundary of the homotopy square traversed by $\eta$ under the isotopy. It proves $d_b\circ\mathrm{push}=\operatorname{conj}$. The filled mapping class gives the same outer automorphism as $d_b$ by construction. Thus the last sequence maps to
\[
1\to\pi\xrightarrow{\operatorname{conj}}\Aut_{\mathrm{lac}}(\pi)
\to\Out_{\mathrm{lac}}(\pi)\to1
\]
by the identity on the end groups. Injectivity and surjectivity of $d_b$ follow by the elementary exact-sequence argument: a kernel element belongs to point-pushing and maps trivially by the injective inner map, and any target element can first be matched in the outer quotient and then in the kernel.

Finally every permutation of old punctures is realized by half-twists supported in small discs along arcs, and an orientation-reversing homeomorphism can be chosen to fix the finite marked set pointwise. Hence the map to $\{\pm1\}\times\mathfrak S_\nu$ is onto; its kernel is precisely the pure orientation-preserving group. Its index is $2\nu!$. This also identifies the subgroups denoted $\mathcal T^{00}$ on scan 188 and proves the compatibility checks left on scan 189.
\end{proof}
\noindent\emph{Manuscript:} \S18, folios 93--94, scans 186--189. On scan 187 the inequality must be $2g+\nu'\ge4$ when $\nu'$ counts punctures before filling; the meaningful alternative $\le4$ contradicts both examples and the stated hyperbolicity requirement. The mathematical argument supports the first alternative without settling the handwriting.

\begin{claim}{A.13}{Proved}{Filling into the two abelian targets}
Exactly one of the drilled and filled groups is abelian in the two cases $(0,3)\to(0,2)$ and $(1,1)\to(1,0)$. In the first, the homomorphism on automorphisms of the underlying groups with peripheral subgroup classes is
\[
 C_2\times C_2=\Out_{\mathrm{lac}}(F_2;[\langle xy\rangle])
 \longrightarrow\Aut(\Z)=C_2,
\]
with kernel of order two. Its nontrivial element has representative $x\mapsto y^{-1}$, $y\mapsto x^{-1}$. Its restriction to the positive-orientation subgroup is an isomorphism. Retaining the two indexed peripheral source maps and their orientation line on the annulus enlarges the target automorphism group to $C_2\times C_2$ and restores the full isomorphism. In the second case the filling map is the isomorphism $\operatorname{GL}_2(\Z)\to\operatorname{GL}_2(\Z)$, with positive subgroup $\operatorname{SL}_2(\Z)$.
\end{claim}
\begin{proof}
A.1 lists the two cases. For a three-punctured sphere the extended mapping class group is $\mathfrak S_3\times C_2$: a homeomorphism fixing the three punctures is isotopic, in each orientation class, to the identity or a fixed reflection, by cutting along three joining arcs and applying the disc isotopy. The stabilizer of the distinguished puncture is $C_2\times C_2$. On $F_2=\langle x,y\rangle$, the old-puncture transposition $x\leftrightarrow y$ and the reflection $x\mapsto x^{-1},y\mapsto y^{-1}$ both induce inversion after imposing $xy=1$. Their product is the displayed nontrivial kernel element. It takes the distinguished loop to its inverse and reverses orientation. Direct free-word and finite-group calculations appear in the attached program.

For an annulus with ordered end sources $\kappa_1(1)=t$ and $\kappa_2(1)=t^{-1}$, an automorphism includes a sign $\epsilon\in\{\pm1\}$ on $T$ and a permutation $s\in\mathfrak S_2$. Compatibility forces the underlying map $t\mapsto t^{\epsilon\operatorname{sgn}(s)}$. Thus there are four automorphisms of this full data, and the map forgetting $s,\epsilon$ is exactly the preceding quotient. This calculation specifies the information lost by retaining only subgroup images.

For the torus, every matrix in $\operatorname{GL}_2(\Z)$ is realized by a linear torus homeomorphism fixing the origin. A based torus homeomorphism inducing the identity matrix is isotopic relative to that point to the identity. Indeed the restriction of evaluation to the translation subgroup $\R^2/\Z^2$ is the identity of the torus, so evaluation is surjective on fundamental groups and the point-pushing kernel is zero. The unbased kernel is the identity component: isotope the images of two primitive curves intersecting once back to those curves, use the identity matrix to remove their remaining twists, and cut along them to a disc, where the Alexander isotopy finishes the argument. Removing the fixed point identifies this mapping class group with the peripheral-preserving outer group of the once-punctured torus by Dehn--Nielsen--Baer. Filling induces abelianization, hence the asserted identity map on matrices. Its orientation sign is the determinant.
\end{proof}
\noindent\emph{Manuscript:} \S18, folios 94--97, scans 189--194. The abelian cases in the margin of scan 194 refer to the filled types; the corresponding drilled groups are both $F_2$.

\begin{claim}{A.14}{Known}{Primitive loops in finite-step solvable quotients}
Let $\Sigma$ be a nonempty set of primes, $F$ a free pro-$\Sigma$ group of finite rank at least two, and $F^{[j]}$ its closed derived series. Put $F_m=F/F^{[m]}$, $m\ge2$. For any element $x$ of a free basis and any positive integer $n$,
\[
 C_{F_m}(x^n)=N_{F_m}(\overline{\langle x\rangle})
 =\overline{\langle x\rangle}\simeq\Z_\Sigma,
 \qquad \Z_\Sigma=\prod_{p\in\Sigma}\Z_p.
\]
Every open subgroup of $F_m$ has trivial centralizer. For $\nu\ge2$, all peripheral loops of the pro-$\Sigma$ group of type $(g,\nu)$ satisfy these formulas in every $F_m$. Their indexed conjugacy classes are pairwise distinct in all hyperbolic types.
\end{claim}
\begin{proof}
The power-centralizer formula is \cite[Theorem 1.5, finite-rank case]{A_Yamaguchi2026}; the source explicitly treats every nonzero integer power, including powers divisible by $2$ or $3$. We checked its finite-rank proof: Proposition 1.3 uses the finite semidirect product $(\Z/N)[G/\langle\bar x\rangle]\rtimes(G\times\langle\bar x\rangle)$ to force the centralizer into $\langle x\rangle F_m^{[m-1]}$; the coefficient at the distinguished coset forces the first $G$ coordinate into $\langle\bar x\rangle$. At $m=2$, the Blanchfield--Lyndon injection of $F'/F''$ into $\Z_\Sigma[[F^{\mathrm{ab}}]]^r$ and injectivity of multiplication by $\bar x^n-1$ remove the intersection with $F_m'$. That injectivity is proved by passing from $C_M$ to $C_{kM}$: invariant coefficients downstairs are divisible by every $\Sigma$-integer $k$, so are zero. The finite-rank induction passes to finite abelian covers where a positive power of $x$ is a free basis element. These are exactly the finite-rank hypotheses used here; no infinite-rank inverse-limit assertion is used.

If $u$ normalizes $\langle x\rangle$, write $uxu^{-1}=x^\lambda$, $\lambda\in\Z_\Sigma^\times$. The primitive abelianization coordinate gives $\lambda=1$, so the normalizer lies in the already computed centralizer. An open subgroup contains positive powers of two distinct free generators; their cyclic centralizers intersect trivially because their abelianization coordinates are independent. This proves the open-centralizer assertion directly.

Each peripheral generator for $\nu\ge2$ belongs to a free basis after eliminating another boundary generator. For $\nu\ge3$ distinct peripheral abelianization lines separate the conjugacy classes. When $\nu=2$ and $g\ge1$, kill $c_1$ and all but one handle pair, mapping onto a free metabelian group on $a,b$. The image of $c_1$ is $1$ and that of $c_2$ is $[a,b]^{-1}\ne1$; the latter has infinite order in the upper unitriangular $3\times3$ group over $\Z_\Sigma$. Hence they cannot be conjugate, or commensurable with finite indices. These maps factor through every $F_m$, $m\ge2$.
\end{proof}
\noindent\emph{Manuscript:} \S13 bis, scans 117--118, and the strengthened finite-step task in Part A(c).

\begin{claim}{A.15}{Proved}{The compact intersection calculation used in descent}
Let $(H_\alpha)$ be a decreasing directed family of closed subgroups of a profinite group $P$, and put $K=\bigcap_\alpha H_\alpha$. Then
\[
 K^{[j]}=\bigcap_\alpha H_\alpha^{[j]}\quad(j\ge0).
\]
For a closed subgroup $L\subset P$ and decreasing closed normal subgroups $N_\alpha$ of $P$, one also has $\bigcap_\alpha LN_\alpha=L(\bigcap_\alpha N_\alpha)$.
\end{claim}
\begin{proof}
For any open normal $V\subset P$, the images of $H_\alpha$ in the finite group $P/V$ eventually stabilize. Their eventual image equals that of $K$: each compatible coset fibre is a decreasing family of nonempty compact sets. Taking a fixed number of derived subgroups commutes with homomorphic images, and taking the inverse limit of all finite images recovers the closed derived subgroup. This proves the first equality. For the second equality, an element $z$ of every $LN_\alpha$ gives nonempty closed sets $L\cap zN_\alpha$; compactness makes their intersection nonempty. A member of that intersection gives the required single factorization of $z$.
\end{proof}

\begin{claim}{A.16}{Proved}{The one-puncture solvable threshold, with a finite certificate}
Let $F$ be the pro-$\Sigma$ completion of $\pi_{g,1}$, $g\ge1$, let $c=\prod_{j=1}^g[a_j,b_j]$ (the inverse peripheral generator), and write $c$ for its image in $F_m=F/F^{[m]}$. For every positive integer $n$,
\[
\begin{array}{ll}
m=2:& C_{F_2}(c^n)=N_{F_2}(\overline{\langle c\rangle})=F_2',\\[2pt]
m\ge3:&C_{F_m}(c^n)=N_{F_m}(\overline{\langle c\rangle})
 =\overline{\langle c\rangle}.
\end{array}
\]
In the first line the inclusion $\overline{\langle c\rangle}\subset F_2'$ is strict. Thus peripheral rigidity fails at the smallest solvable depth $m=2$ and is restored at every $m\ge3$, for the full profinite and every pro-$p$ version.
\end{claim}
\begin{proof}
Let $H\subset F$ run over open normal subgroups containing $F'$, with $|F/H|\ge3$. They form a family whose intersection is $F'$. The corresponding finite abelian covering of the once-punctured surface has exactly $d=|F/H|$ punctures, because $c\in F'$ has trivial image in the deck group. Choose the cusp above the original based peripheral loop. Its loop is $c$ itself; the other loops are its conjugates by deck representatives. The submodule of $H^{\mathrm{ab}}$ generated by these cusps is
\[
 \Z_\Sigma^{F/H}\big/\Z_\Sigma\!\left(\sum_{q\in F/H}e_q\right),
\]
with the deck group acting by left translation. This follows either from the relative surface homology sequence of A.2 or directly from the free-group peripheral presentation. For $d\ge3$ the lines generated by distinct $e_q$ in this module are distinct; moreover $ne_q=ae_1$, for $0\ne n\in\Z$ and $a\in\Z_\Sigma$, forces $q=1$ and $a=n$. To check this last assertion, lift the equality to the free module and compare a coordinate different from both $1$ and $q$; its coefficient forces the possible multiple of $\sum e_q$ to be zero, after which the coordinates at $q$ and $1$ give the result.

For $m\ge2$, $F^{[m]}\subset H'$, so the natural map $H\to H^{\mathrm{ab}}$ factors through the image of $H$ in $F_m$. Let $u\in F$ represent an element of $C_{F_m}(c^n)$. The relation $uc^nu^{-1}=c^n$ modulo $F^{[m]}$ implies $ne_{\bar u}=ne_1$ in the cusp submodule of $H^{\mathrm{ab}}$. Hence $u\in H$ for every such $H$, so its image belongs to $F_m'$. The same argument for a normalizer, using $uc u^{-1}=c^\lambda$ modulo $F^{[m]}$, $\lambda\in\Z_\Sigma^\times$, gives $u\in F'$ and $\lambda=1$. The subgroup $\overline{\langle c\rangle}$ is indeed $\Z_\Sigma$: the unitriangular $3\times3$ quotient used in A.1 proves injectivity of all its prime components. Thus its normalizer equals its centralizer.

When $m=2$, the derived subgroup is abelian, contains $c$, and centralizes it. This gives equality with $F_2'$. Now let $m\ge3$ and let $uF^{[m]}$ centralize $c^n$. We have proved $u\in F'\subset H$. Since $F'\subset H$,
$F^{[m]}=(F')^{[m-1]}\subset H^{[m-1]}$; therefore there is a well-defined continuous epimorphism
\[
 H/F^{[m]}\longrightarrow H/H^{[m-1]}.
\]
The group $H$ is free pro-$\Sigma$, and $c$ is a member of a free basis because this cover has at least two punctures. By A.14, applied at depth $m-1\ge2$, the image of $u$ in this quotient belongs to $\overline{\langle c\rangle}$. Consequently
\[
 u\in\bigcap_H\overline{\langle c\rangle}H^{[m-1]}
 =\overline{\langle c\rangle}\left(\bigcap_H H^{[m-1]}\right)
 =\overline{\langle c\rangle}(F')^{[m-1]}
 =\overline{\langle c\rangle}F^{[m]},
\]
where both interchanges are proved in A.15. This proves the second line.

For strictness in the first line, set $a=a_1$, $b=b_1$, and kill the remaining handle pairs. In $\operatorname{UT}_4(\F_p)$ take
\[
 x=1+E_{12}+E_{23},\qquad y=1+E_{34}.
\]
Their group is metabelian; direct multiplication gives $[x,y]=1+E_{24}$ and $x[x,y]x^{-1}=1+E_{24}+E_{14}$. The second matrix is not a power of the first. Thus $aca^{-1}\in F_2'$ centralizes $c$ but does not belong to its procyclic closure. For every $p\in\Sigma$ this quotient factors through $F_2$ and proves strictness; the attached program independently enumerates and verifies the cases $p=2,3,5$. In particular there is no excluded residue-characteristic caveat.
\end{proof}
\noindent\emph{Manuscript:} \S13 bis, scan 118, in the strengthened finite-step question of Part A(c). This assertion does not assert a new literature priority; the searched sources and the proof audit are listed in the report.

\begin{claim}{A.17}{Proved}{Centre and torsion in all finite-step surface quotients}
Let $P$ be the full profinite or pro-$\Sigma$ completion of a hyperbolic surface group, open or closed, and let $P_m=P/P^{[m]}$. For every $m\ge2$, $P_m$ is torsion-free and centre-free. This verifies the form of \cite[Theorem 2.9 and Corollary 2.10]{A_Yamaguchi2026} needed here by the following argument from finite covers and transfer.
\end{claim}
\begin{proof}
Every open subgroup of $P$ is the pro-$\Sigma$ group of a finite covering surface, so its abelianization is a free $\Z_\Sigma$-module. A closed subgroup $K$ is the intersection of its open neighbourhood subgroups; by A.15 its abelianization is the inverse limit of their abelianizations. More explicitly, a compatible system in those abelianizations lifts to a compatible coset system in the compact group, and the kernel is the intersection of the closed derived subgroups. Thus $K^{\mathrm{ab}}$ is torsion-free. Applying this to each $P^{[j]}$ and inducting in
\[
1\to(P^{[m-1]})^{\mathrm{ab}}\to P_m\to P_{m-1}\to1
\]
proves torsion-freeness, starting with the torsion-free abelianization.

We also need the faithful action of a finite deck group $Q=H/N$ on $N^{\mathrm{ab}}$, for open $N\trianglelefteq H\subset P$. Choose $\ell\in\Sigma$. If a subgroup $Q_0$ acts trivially on $H^1$ of the covering surface with $\Q_\ell$ coefficients, transfer gives an isomorphism of that $H^1$ with the $H^1$ of the quotient by $Q_0$. The two surfaces are both proper or both nonproper, and their Euler characteristics are consequently equal. Multiplicativity of Euler characteristic gives $\chi(Y)=|Q_0|\chi(Y/Q_0)$. Both are negative, so $Q_0=1$. The comparison with $N^{\mathrm{ab}}\otimes_{\Z_\Sigma}\Q_\ell$ is the canonical duality between first homology and first cohomology. Thus the deck action on $N^{\mathrm{ab}}$ is faithful.

Put $K=P^{[m-1]}$, $A=K/K'$, and $Q=P/K'$. For an open normal subgroup $N\subset Q$ containing $A$, its preimage $\widetilde N\subset P$ contains $K$, hence $K'\subset\widetilde N'$. Therefore $N^{\mathrm{ab}}=\widetilde N^{\mathrm{ab}}$, and the preceding deck faithfulness implies that an element of $Z(Q)$ belongs to every such $N$. Their intersection is $A$.

We finish by proving $A^{P/K}=0$. For any open normal $V\subset P$ containing $K$, write $j:V^{\mathrm{ab}}\to P^{\mathrm{ab}}$ for inclusion and $t:P^{\mathrm{ab}}\to V^{\mathrm{ab}}$ for transfer. The composite $tj$ on the $P/V$-invariants is multiplication by $[P:V]$, which is injective because $V^{\mathrm{ab}}$ is torsion-free. Hence inclusion is injective on those invariants. Passing to the inverse limit over $V$ shows that $K^{\mathrm{ab},P/K}\to P^{\mathrm{ab}}$ is injective. But $K\subset P'$, so this map is zero. Therefore $A^{P/K}=0$. Since $Z(Q)\subset A$, the centre is zero, as claimed. This proof checks precisely the transfer and finite-cover steps used in the cited unrefereed result.
\end{proof}

\begin{claim}{A.18}{Proved}{The remaining affine solvable rigidity list and a closed bound}
For all hyperbolic types with $\nu>0$, and all $m\ge2$, the maximal $m$-step-solvable full profinite or pro-$\Sigma$ quotient has trivial centre and trivial centralizers of open subgroups, and is torsion-free. Peripheral conjugacy classes remain distinct. For $\nu\ge2$ all peripheral normalizer and integer-power-centralizer equalities hold at every $m\ge2$; for $\nu=1$ their exact values are those of A.16. In an extension obtained by quotienting a geometric action extension by the $m$th closed derived subgroup of its kernel, the equality $N_E(L_s)=E_s$ holds in the former cases and at $m\ge3$ in the latter. At $\nu=1,m=2$ its correct value is $N_E(L_s)=F_2'E_s$, with the intersection $F_2'\cap E_s=L_s$. In all these cases every transporter $\operatorname{Trans}_E(L_s^d,L_s)$, $d\ge1$, equals the same displayed normalizer. Since the kernels are centre-free by A.17, the extension reconstruction of A.6 also applies to all these solvable quotients. Standard free-basis generators satisfy the power-centralizer formula at every $m\ge2$ by A.14.
For a closed hyperbolic group $P$ and an open subgroup $H\subset P_m$, the proven universal bound is
\[
 C_{P_m}(H)\subseteq P_m^{[m-1]}.
\]
\end{claim}
\begin{proof}
The first assertions are A.14, A.16, and A.17. For the extension calculation, the quotient map is injective on the peripheral group, since the peripheral generator has infinite order in $P_2$ by A.1. The image of a geometric peripheral lift group still maps onto the puncture stabilizer. The kernel of the normalizer map to that stabilizer is the normalizer calculated in A.14 or A.16, so the same lifting argument as A.5 proves the equalities. In the metabelian one-puncture case there is just one puncture, so the stabilizer is the whole acting group and multiplication by $E_s$ gives the asserted $F_2'E_s$.

It remains to justify the one-sided transporter equality in these quotients. Suppose $eL_i^de^{-1}\subset L_i$. The action sends $L_i$ to a conjugate of some $L_j$ and its chosen generator to a conjugate of $c_j^\lambda$, with $\lambda\in\Z_\Sigma^\times$. For $\nu\ge3$, the primitive abelianization lines imply $j=i$. For $\nu=2$, $g\ge1$, the quotient killing $c_i$ and all but one handle pair sends $c_j$ to a nonzero commutator of infinite order in $\operatorname{UT}_3(\Z_\Sigma)$; thus $j\ne i$ would contradict the given containment. Choose an element of $E_i$ with the same image as $e$ in the acting group and divide it out. The resulting element $u$ lies in $P_m$ and satisfies $uc_i^du^{-1}=c_i^a$ for some $a\in\Z_\Sigma$. The primitive abelianization coordinate forces $a=d$, so A.14 puts $u$ in $L_i$. This proves the equality for $\nu\ge2$. For $\nu=1$, the acting group already fixes the only puncture. After the same division by $E_i$, apply the finite abelian cover calculation in A.16 to $uc^du^{-1}=c^a$. It forces $a=d$ and reduces the transporter to the integer-power centralizer calculated there. This gives $E_i$ at $m\ge3$ and $F_2'E_i$ at $m=2$.

The closed bound recorded here is strengthened to triviality in A.19 below. For this bound, put $A=P_m^{[m-1]}$ and choose any open normal $N\subset P_m$ containing $A$. The group $N^{\mathrm{ab}}$ is the torsion-free abelianization of its preimage in $P$, as in A.17. If $z$ centralizes $H$, its conjugation action fixes the image of $H\cap N$ in $N^{\mathrm{ab}}$. This image has finite index, so some nonzero integer multiple of every element is fixed. Torsion-freeness then makes the action trivial on all of $N^{\mathrm{ab}}$. Deck faithfulness from A.17 implies $z\in N$. Intersecting all these $N$ yields $z\in A$.
\end{proof}
\noindent\emph{Manuscript:} \S13 bis, scan 118; \S35, scans 498--501; additional finite-step task A(c).

\begin{claim}{A.19}{Proved}{Closed finite-step power centralizers and slimness}
Let $P$ be the pro-$\Sigma$ completion of a closed oriented surface group of genus $g\ge2$, let $m\ge2$, and put $P_m=P/P^{[m]}$. Let $x$ be represented by a nonseparating simple closed curve, in particular by any standard handle generator. Then for every positive integer $n$,
\[
 C_{P_m}(x^n)=N_{P_m}(\overline{\langle x^n\rangle})
 =\overline{\langle x\rangle}\simeq\Z_\Sigma.
\]
Consequently $C_{P_m}(H)=1$ for every open subgroup $H$. This includes the full profinite and every pro-$p$ quotient, without a restriction at $p=2$ or $3$.
\end{claim}
\begin{proof}
We first give the completed module construction required for the metabelian case. Let $F$ be a free pro-$\Sigma$ group on $z_1,\ldots,z_r$, let $A=F/F'$, and put $R=\Z_\Sigma[[A]]$. The continuous homomorphism
\[
 \mu:F\longrightarrow R^r\rtimes A,
 \qquad z_i\longmapsto(e_i,\bar z_i),
\]
exists by freeness; $A$ acts by multiplication on the coefficient module, so the product in the target is $(v,a)(w,b)=(v+aw,ab)$. Its kernel is exactly $F''$. The inclusion $F''\subset\ker\mu$ follows because the target is metabelian. To prove the reverse inclusion without assuming any exactness result about completed chain complexes, take an arbitrary finite metabelian quotient $G$ of $F$. Write $K=G'$, $Q=G/K$, and choose a section $t:Q\to G$ with $t(1)=1$. There is an injective homomorphism
\[
 G\longrightarrow K^Q\rtimes Q,
 \qquad g\longmapsto(f_g,\bar g),\qquad
 f_g(q)=t(q)g\,t(q\bar g)^{-1}.
\]
Indeed $f_{gh}(q)=f_g(q)f_h(q\bar g)$, and, when $\bar g=1$, evaluation at $q=1$ recovers $g$. Use additive notation for the finite abelian group $K$. The finite module $K^Q$ is an $R$-module through $A\to Q$ and right translation on $Q$. Sending $e_i$ to the function coordinate of the image of $z_i$ gives a continuous $R$-linear map $R^r\to K^Q$. The resulting homomorphism of semidirect products makes the homomorphism from $F$ factor through $\mu$, since it agrees on every free generator. Hence $\ker\mu$ maps trivially to every finite metabelian quotient of $F$. These quotients separate $F/F''$, proving the assertion. In particular the function coordinate gives an injective $R$-linear morphism
\[
 F'/F''\lhook\joinrel\longrightarrow R^r.
\]
Its image is a closed $R$-submodule: conjugation gives the $A$-action, integral powers give the scalar action, and continuity extends these to the completed group ring.

Apply this construction with $r=2g$ and
\[
 P=F/\langle\!\langle r_0\rangle\!\rangle,
 \qquad r_0=\prod_{i=1}^g[a_i,b_i],
 \qquad A=P^{\mathrm{ab}}=\Z_\Sigma^{2g}.
\]
The function coordinate of $r_0$ is
\[
 v=(1-\bar b_1,\bar a_1-1,\ldots,1-\bar b_g,\bar a_g-1)\in R^{2g}.
\]
This is obtained by multiplying the four factors of each commutator in $R^{2g}\rtimes A$. Since $r_0\in F'$, its closed normal closure maps in $F'/F''$ onto exactly $Rv$: its conjugates give the $A$-multiples of $v$, their products give their integral sums, and their closure is $Rv$. The last module is closed because it is the image of the compact ring $R$. Thus the natural quotient maps induce the injective comparison
\[
 B:=P'/P''\lhook\joinrel\longrightarrow R^{2g}/Rv.                 \tag{A.19.1}
\]
Indeed $P''$ is the image of $F''$ and the kernel in $F'/F''$ is precisely the just calculated module. This establishes the needed completed Fox-module comparison directly.

Choose a geometric basis with $x=a_1$, and set $h=\bar a_1^n-1$. We use two elementary facts about these completed group rings. First, multiplication by $h$ on $R$ is injective. Second, multiplication by $\bar a_2-1$ on $R/hR$ is injective. Here
\[
 R/hR\xrightarrow{\sim}
 \Z_\Sigma[[A/\overline{\langle\bar a_1^n\rangle}]]                \tag{A.19.2}
\]
is the map induced by the group quotient. Its kernel is the closed ideal generated by $h$; that ideal equals $hR$ since $hR$ is compact. To prove the injectivity assertions, one can use the following coefficient calculation. In a completed group ring of $\Z_\Sigma\times D$, consider an element invariant under the chosen infinite cyclic coordinate, or under its $n$th power. At finite coefficient level $\Z/k$ and finite coordinate level $C_M$, enlarge the coordinate level to $C_{kM}$, taking $M$ divisible by the $\Sigma$-part of $n$ in the second case. Invariance makes the coefficients constant on each orbit. The projection down to $C_M$ sums $k$ equal coefficients on each orbit fibre and therefore is zero modulo $k$. This works with every finite quotient of the remaining factor $D$. All finite images of the original element vanish, so the element is zero. Apply this argument first to the $a_1$ coordinate of $R$, and then to the independent $a_2$ coordinate in the group of (A.19.2). It applies to arbitrary $\Sigma$, because $k$ and the chosen coordinate orders have prime factors in $\Sigma$.

These ring facts show that multiplication by $h$ is injective on $R^{2g}/Rv$. If $hz=vt$, inspect the $b_2$ coordinate, which is $h z_{b_2}=(\bar a_2-1)t$. Reducing modulo $h$ and using the second injectivity gives $t=ht'$ for some $t'\in R$. The first injectivity then gives $z=vt'$ in every coordinate. Consequently (A.19.1) implies
\[
 B^{\langle x^n\rangle}=0.                                      \tag{A.19.3}
\]

For completeness, we next constrain the abelianization of an arbitrary centralizer element, at every depth $m\ge2$. Put $V=\Z_\Sigma^{2g}$ with its ordered handle basis and put $W=\bigwedge^2V/\Z_\Sigma\omega$, where $\omega=\sum_i a_i\wedge b_i$. For every $\Sigma$-integer $k$, form the finite class-two group with underlying set $(V/kV)\times(W/kW)$ and product
\[
 (q,w)(q',w')=(q+q',w+w'+\beta(q,q')),
 \qquad
 \beta(q,q')=\sum_{i<j}q_iq'_j\,z_i\wedge z_j.
\]
Its commutator is $(0,q\wedge q')$. The standard surface relation maps to $\omega=0$, so sending the handle generators to their coordinate vectors defines a quotient homomorphism from $P$. The finite group is metabelian and a $\Sigma$-group; the homomorphism therefore factors through every $P_m$, $m\ge2$. If $u\in P_m$ centralizes $a_1^n$ and has abelianization $q$, these homomorphisms give
$n(q\wedge a_1)=0$ in $W$. The module $W$ is free because $\omega$ is primitive, so multiplication by the nonzero integer $n$ is injective. Write $q\wedge a_1=t\omega$. The coefficient of $a_2\wedge b_2$ on the left is zero; on the right it is $t$. Hence $t=0$, and the remaining coordinates give $q=\lambda a_1$ for some $\lambda\in\Z_\Sigma$. The procyclic group generated by $a_1$ injects into $V$, so dividing $u$ by $a_1^\lambda$ puts it in $P_m'$ without changing its centralization property. At $m=2$, (A.19.3) makes this remaining element trivial, proving the centralizer formula.

We now induct on $m$, simultaneously for every genus at least two and every nonseparating simple closed curve. Suppose $m\ge3$, and divide a centralizer element as just explained so that it is represented by $u\in P'$. Let $H$ run over all open normal subgroups of $P$ containing $P'$, and let $d$ be the order of $x$ in $P/H$. The corresponding finite abelian covering of the closed surface has genus at least two. The based lift of the simple curve $x$ is represented by $x^d$, traversed once. This lift is a simple nonseparating curve: its integral homology class pushes forward to $d[x]\ne0$, whereas a separating curve has zero homology class. It is therefore a standard handle generator for some geometric basis of the covering surface group $H$.

Since $P'\subset H$, we have $P^{[m]}\subset H^{[m-1]}$. The quotient morphism
\[
 H/P^{[m]}\longrightarrow H/H^{[m-1]}
\]
is consequently defined. The image of $u$ centralizes $(x^d)^n$ in its target. By the induction hypothesis at depth $m-1$, that image belongs to $\overline{\langle x^d\rangle}$. Hence
\[
 u\in\bigcap_H\overline{\langle x\rangle}H^{[m-1]}
 =\overline{\langle x\rangle}\left(\bigcap_HH\right)^{[m-1]}
 =\overline{\langle x\rangle}P^{[m]},
\]
by A.15 and $\bigcap_HH=P'$. The intersection of $\overline{\langle x\rangle}$ with $P'$ is trivial, again by its primitive abelianization coordinate. Since $u\in P'$, its image in $P_m$ is trivial. This completes the induction.

If an element normalizes $\overline{\langle x^n\rangle}$, its conjugation sends $x^n$ to $x^{n\lambda}$ for a unit $\lambda\in\Z_\Sigma^\times$. Abelianization gives $n\lambda=n$, whence $\lambda=1$, so the normalizer equals the calculated centralizer. Finally an open subgroup contains positive integer powers of $a_1$ and $b_1$. Their calculated centralizers have trivial intersection because they inject into distinct independent abelianization coordinates. This proves slimness.
\end{proof}
\noindent\emph{Manuscript:} \S35, scan 501, equation (4), and the finite-step strengthening in Part A(c). The proof establishes the exact module comparison and all closed-group descent maps used; it does not rely on the unknown general slimness statement mentioned before Proposition 2.6 in \cite{A_Yamaguchi2026}, nor claim a literature priority.

\begin{claim}{A.20}{Proved}{The typed completion comparison for filling}
For the two hyperbolic groups in A.12, put $\Pi'=\widehat{\pi'}$ and $\Pi=\widehat\pi$. There is a commutative diagram of profinite groups
\[
\begin{array}{ccc}
 \widehat{\Out_{\mathrm{lac}}(\pi';b)}
 &\xrightarrow{\ \widehat{d_b}\ }&
 \widehat{\Aut_{\mathrm{lac}}(\pi)}\\
 \big\downarrow c_{\mathrm{out}}&&\big\downarrow c_{\mathrm{aut}}\\
 \Out_{\mathrm{lac}}(\Pi';b)
 &\xrightarrow{\ d_b^{\mathrm{cont}}\ }&
 \Aut_{\mathrm{lac}}(\Pi).
\end{array}
\]
The upper horizontal map is an isomorphism. The lower horizontal map is the normalized filling homomorphism on all continuous loop automorphisms. The vertical maps are induced by extending discrete automorphisms to group completions. The same construction with pro-$p$ completions of the surface groups still has the full profinite completions of the discrete automorphism groups in its upper row. These assertions specify the precise consequence of the discrete filling theorem; no isomorphism assertion for the lower horizontal or vertical arrows is used.
\end{claim}
\begin{proof}
The presentation gives $\Pi'/\overline{\langle\!\langle L_b\rangle\!\rangle}=\Pi$. By A.7, $N_{\Pi'}(L_b)=L_b$. Normalize an outer representative to preserve $L_b$, then pass to this quotient. Two normalized representatives differ by conjugation by an element of $L_b$, which dies in the quotient. This proves that $d_b^{\mathrm{cont}}$ is a well-defined homomorphism. It is continuous: the closed subgroup of $\Aut(\Pi')$ preserving $L_b$ is compact, its quotient onto the indicated outer group has kernel the inner automorphisms by $L_b$, and the continuous quotient action factors through this quotient. The same proof works for the pro-$p$ groups using A.7.

Every discrete automorphism extends uniquely to an automorphism of its completion; extensions of inner automorphisms remain inner. Since the resulting target automorphism groups are profinite, the extension maps themselves extend to the full profinite completions of the discrete automorphism groups. Normalizing a discrete representative and then completing gives the same quotient automorphism as completing and then normalizing. Thus the square commutes on the dense discrete subgroup and therefore everywhere. A.12 implies that its upper arrow is an isomorphism. The quotient maps $\widehat\pi\twoheadrightarrow\pi^{(p)}$ and $\widehat{\pi'}\twoheadrightarrow(\pi')^{(p)}$ have characteristic kernels and commute with filling, giving the further comparison from the full square to its pro-$p$ version.

For odd $p$, replacing the upper groups by their pro-$p$ completions would in general fail to define these same maps: an orientation-reversing involution is killed by every homomorphism to a pro-$p$ group, whereas its action on the pro-$p$ orientation line is multiplication by $-1\ne1$. This is an explicit obstruction to that proposed factorization, rather than an identification of two differently typed completion operations.
\end{proof}
